\section{Further examples and relative tangent slopes}\label{sec:auxiliary}

We record several consequences of the methods used above. All slopes
in this appendix are taken with respect to the anticanonical
polarization of the total space.

\subsection{Del Pezzo and pseudo del Pezzo varieties}

Let $n=2m\ge2$, let $e_1,\ldots,e_n$ be the standard basis of
$\ZZ^n$, and put $s=e_1+\cdots+e_n$. The del Pezzo variety $V_n$ and
the pseudo del Pezzo variety $\widetilde V_n$ have ray generators
\[
 V_n:\quad \pm e_1,\ldots,\pm e_n,\pm s,
 \qquad
 \widetilde V_n:\quad \pm e_1,\ldots,\pm e_n,s,
\]
respectively; their fans consist of the cones over the faces of the
convex hulls of these generators
\cite{VoskresenskiiKlyachko1985,Ewald1988,Nill2006}.
Thus $V_2$ is the del Pezzo surface of degree six and
$\widetilde V_2$ the one of degree seven.

\begin{proposition}\label{aux:delpezzo}
For every $m\ge1$, both $V_{2m}$ and $\widetilde V_{2m}$ are smooth
toric Fano varieties with stable tangent bundles. Their Picard numbers
are $2m+2$ and $2m+1$, respectively.
\end{proposition}

\begin{proof}
Their anticanonical polytopes are
\[
 [-1,1]^n\cap\left\{\left|\sum_i x_i\right|\le1\right\},
 \qquad
 [-1,1]^n\cap\left\{\sum_i x_i\ge-1\right\}.
\]
Here $n$ is even, so no cube vertex lies on a cutting hyperplane.
The remaining cube vertices have coordinate sum zero in the first
case and nonnegative coordinate sum in the second. Every new vertex
lies on a cube edge: one coordinate is zero and the remaining
coordinates are $\pm1$ with sum $1$ or $-1$ in the first case, and
sum $-1$ in the second. The primitive normals of the active facets
are a signed coordinate basis at an old vertex. At a new vertex they
are $n-1$ signed coordinate vectors together with $s$ or $-s$; their
determinant is $\pm1$. Both polytopes are therefore integral and
simple with unimodular normal cones. The origin is interior, and
all displayed inequalities define facets. This proves the smooth
Fano assertion.

For $V_n$, use the equivalent lattice presentation
$\ZZ^{n+1}/\ZZ(e_0+\cdots+e_n)$, with rays $\pm e_i$.
The symmetric group $S_{n+1}$ acts through its irreducible standard
representation on $N_\CC$. Its ray matroid is connected: opposite
rays are parallel and $\{e_0,\ldots,e_n\}$ is a circuit meeting
every parallel class. Proposition~\ref{prop:symmetry} proves stability.

For $\widetilde V_n$, the group $S_n$ permuting the coordinate
vectors preserves the fan. Its only nonzero proper invariant
subspaces are $\CC s$ and $\{\sum_i x_i=0\}$; this follows from the
decomposition into the trivial and standard representations,
which are non-isomorphic irreducibles, including when $n=2$.
The latter hyperplane contains no ray. The circuit
$\{e_1,\ldots,e_n,s\}$ and the opposite pairs show that the ray
matroid is connected. Write $d_{e_i}=p$, $d_{-e_i}=q$, and $d_s=h$.
Lemma~\ref{lem:degreerelation} gives $q=p+h$, so
\[
 D=n(p+q)+h=n(2p+h)+h>nh.
\]
Thus the only invariant subspace spanned by rays, $\CC s$, satisfies
the strict slope inequality. Proposition~\ref{prop:symmetry} applies.
Finally, subtracting the dimension from the number of rays gives
the two Picard numbers.
\end{proof}

\subsection{Extensions with base \texorpdfstring{$\PP^1$}{P1}}

Let $Z=X_{\Sigma_Z}$ be a positive-dimensional smooth projective toric variety with lattice
$N_Z$, and let $t\in N_Z$. In $N_Z\oplus\ZZ$, take the cones
generated by a cone of $\Sigma_Z$ and either $(t,1)$ or $(0,-1)$,
together with all their faces. Denote the resulting variety by $X_t$;
projection to the last coordinate gives a toric bundle
$X_t\to\PP^1$ with fibre $Z$.

\begin{proposition}\label{aux:baseextension}
The fan of $X_t$ is smooth and complete. The variety $X_t$ is Fano
if and only if $Z$ is Fano and
\[
 t\in Q\cap N_Z,\qquad
 Q=\operatorname{conv}\{u_\rho:\rho\in\Sigma_Z(1)\}.
\]
When $Z$ is Fano, $Q\cap N_Z$ consists exactly of the origin and
the vertices of $Q$.
\end{proposition}

\begin{proof}
On the upper half-space, the shear $(x,y)\mapsto(x-yt,y)$ identifies
the cones with the product of $\Sigma_Z$ and the nonnegative ray.
On the lower half-space they already have product form. This proves
completeness, and each maximal cone is unimodular. If $X_t$ is Fano,
restriction of $-K_{X_t}$ to a fibre gives $-K_Z$, so $Z$ is Fano.

Assume henceforth that $Z$ is Fano. For a facet $F$ of $Q$, let
$u_F\in M_Z$ satisfy $\langle u_F,v\rangle=1$ on $F$ and be
strictly less than one at every other vertex. The linear forms
supporting the two proposed facets over $F$ are
\[
 (u_F,-1),\qquad (u_F,1-\langle u_F,t\rangle).
\]
Evaluating each on the opposite apex shows that the strict convexity
condition is $\langle u_F,t\rangle<2$. Since this quantity is integral,
the condition is equivalent to $\langle u_F,t\rangle\le1$ for every
facet, hence to $t\in Q$. These are exactly the strict support-function
inequalities for ampleness of $-K_{X_t}$. Finally, the simplices
$\operatorname{conv}(0,F)$ cover $Q$ and are unimodular, so they
contain no lattice points other than their vertices.
\end{proof}

\subsection{Bundles over products of projective lines}

Let $F$ be a smooth toric Fano variety of dimension $f\ge1$, and set
\[
 P_F=\{m\in M_{F,\RR}:\langle m,v\rangle\ge-1
                       \text{ for every ray generator }v\text{ of }F\}.
\]
For $k\ge1$, choose each $t_i$ to be zero or a primitive ray
generator of $F$. The bundle $X=X_F(t_1,\ldots,t_k)$ over $(\PP^1)^k$
has rays $(v,0)$, $(t_i,e_i)$, and $(0,-e_i)$, with maximal cones
obtained by choosing a maximal fibre cone and one ray from each
base pair. Proposition~\ref{aux:baseextension}, applied successively,
shows that $X$ is smooth Fano. Put $n=f+k$, $[k]=\{1,\ldots,k\}$, and
\[
 I(S)=\int_{P_F}\prod_{i\in S}\langle m,t_i\rangle\,dm,
 \qquad S\subseteq\{1,\ldots,k\},
\]
where $dm$ gives the character lattice covolume one and an empty
product is one.

\begin{proposition}\label{aux:prism}
The anticanonical degree is
\[
 D=n!\int_{P_F}\prod_{i=1}^k(2+\langle m,t_i\rangle)\,dm
   =n!\sum_{S\subseteq[k]}2^{k-|S|}I(S).
\]
For each $i$, the two lifted base rays have equal divisor degree
$(n-1)!J_i$, where
\[
 J_i=\int_{P_F}\prod_{j\ne i}(2+\langle m,t_j\rangle)\,dm.
\]
The relative tangent bundle satisfies
\begin{equation}\label{aux:relativeprism}
 \mu(T_{X/(\PP^1)^k})\ge\mu(T_X)
 \quad\Longleftrightarrow\quad
 \sum_{S\subseteq[k]}|S|\,2^{k-|S|}I(S)\ge0.
\end{equation}
The same equivalence holds with strict inequalities.
\end{proposition}

\begin{proof}
The anticanonical polytope is
\[
 \{(m,y):m\in P_F,\ -1-\langle m,t_i\rangle\le y_i\le1
                              \text{ for }1\le i\le k\}.
\]
All interval lengths are positive. Integration gives $D=n!\operatorname{vol}(P_X)$.
Each base facet projects by a lattice isomorphism onto the region
with the $i$-th interval omitted; this gives its degree $(n-1)!J_i$.
The fibre subspace corresponds to the relative tangent bundle,
which has rank $f$ and degree $D-2(n-1)!\sum_iJ_i$.
Comparing its slope with $D/n$ is equivalent to
$kD\ge2n!\sum_iJ_i$. A term $I(S)$ occurs in the last sum exactly
$k-|S|$ times, with coefficient $2^{k-1-|S|}$ each time.
Subtracting gives \eqref{aux:relativeprism}.
\end{proof}

For $k=1$, write $b_F=\operatorname{vol}(P_F)^{-1}\int_{P_F}m\,dm$.
The relative tangent slope is at least $\mu(T_X)$ exactly when
$\langle b_F,t_1\rangle\ge0$. Thus a Fano toric bundle of this form
over $\PP^1$ cannot have stable tangent bundle if $b_F=0$.

For $k=2$ and centrally symmetric $P_F$, the condition is
\[
 B(t_1,t_2):=\int_{P_F}\langle m,t_1\rangle\langle m,t_2\rangle\,dm\ge0.
\]
For $F=V_{2m}$, use its quotient-lattice presentation with rays
$\pm e_0,\ldots,\pm e_{2m}$. The form $B$ is positive definite and
$S_{2m+1}$-invariant. Irreducibility of the standard representation
implies that it is a positive multiple of the form with entries
$\delta_{ij}-1/(2m+1)$. Hence, for nonzero twists
$t_1=\varepsilon_1e_i$, $t_2=\varepsilon_2e_j$, with
$\varepsilon_1,\varepsilon_2\in\{1,-1\}$, the relative tangent
bundle has slope greater than $\mu(T_X)$ precisely when
\[
 t_2=t_1\quad\text{or}\quad
 i\ne j\ \text{ and }\ \varepsilon_1\varepsilon_2=-1.
\]
For opposite twists with the same index, and for equal signs with
different indices, this particular subsheaf has smaller slope.
These comparisons alone do not decide stability of $T_X$.

\subsection{Projective-line bundles over a toric base}

Let $Y$ be a smooth projective toric variety of dimension $n-1\ge1$, and
let $X=\PP_Y(\mathcal O\oplus L)$ be a toric projective-line bundle.
Choose its fan presentation with fibre rays $(0,\pm1)$ and lifted
base rays $(u_\rho,a_\rho)$. Set $L'=\sum_\rho a_\rho D_\rho$ on $Y$;
the line bundle $L$ corresponds to $L'$ or $-L'$, according to the
projectivization convention. Define
\[
 g(y)=(-K_Y+yL')^{n-1}.
\]

\begin{proposition}\label{aux:p1bundle}
Assume $X$ is Fano. Then $-K_Y\pm L'$ are ample, and
\[
 D=n\int_{-1}^1g(y)\,dy,\qquad
 d_{(0,1)}=g(-1),\qquad d_{(0,-1)}=g(1).
\]
Consequently
\begin{align*}
 \mu(T_{X/Y})\ge\mu(T_X)
 &\Longleftrightarrow
 g(1)+g(-1)\ge\int_{-1}^1g(y)\,dy\\
 &\Longleftrightarrow
 \sum_{\substack{2\le j\le n-1\\ j\text{ even}}}
 \frac{j}{j+1}\binom{n-1}{j}(-K_Y)^{n-1-j}\cdot(L')^j\ge0.
\end{align*}
For $n=3$ this is $(L')^2\ge0$, and for $n=4$ it is
$(-K_Y)\cdot(L')^2\ge0$.
\end{proposition}

\begin{proof}
The two sections have normal bundles $L$ and $L^{-1}$; adjunction
therefore identifies the restrictions of $-K_X$ with $-K_Y+L'$ and
$-K_Y-L'$. They are ample, hence in particular nef. Their convex
combinations $-K_Y+yL'$ are ample for $-1\le y\le1$.
At height $y$, the anticanonical polytope has slice
\[
 P_y=\{m:\langle m,u_\rho\rangle\ge-1-a_\rho y
                         \text{ for every }\rho\},
\]
the polytope of $-K_Y+yL'$, of volume $g(y)/(n-1)!$.
Integrating gives $D$; the facets at $y=-1$ and $y=1$ give the
two divisor degrees. The relative tangent line bundle is
$\mathcal O_X(D_{(0,1)}+D_{(0,-1)})$, so its degree is their sum.
Finally, expand $g(y)=\sum_j c_jy^j$. The difference between the
endpoint sum and the integral is
$2\sum_{j\ge2,\ j\text{ even}}c_jj/(j+1)$, giving the assertion.
\end{proof}

\begin{example}\label{aux:p1counterexample}
A smaller relative tangent slope does not imply stability. Let
$Y=\PP^2\times\PP^1$ and $L'=2H-F$, where $H$ and $F$ are the
hyperplane classes of the two factors. Both $-K_Y\pm L'$ are ample;
the strict support-function inequalities on the two sets of maximal
bundle cones then show that $X$ is Fano. Its rays may be written
\[
 -e_1,\ -e_4,\ e_2-e_4,\ -e_2,\ e_4,\ -e_3,\ e_1+e_3+2e_4.
\]
This is the variety F.4D.0019 in the smooth Fano classification
\cite{polyDB}.
Here $g(y)=3(3+2y)^2(2-y)$, so
\[
 D=400,\qquad\mu(T_X)=100,\qquad\deg T_{X/Y}=g(1)+g(-1)=84.
\]
The divisor degrees in the displayed ray order are
$(64,75,62,62,9,64,64)$. For example, the facet for $-e_1$ has degree
$6\int_{-1}^1(3+2y)(2-y)\,dy=64$, while the two facets for the
$\PP^1$ base directions have degree
$3\int_{-1}^1(3+2y)^2\,dy=62$.
The plane $\operatorname{span}(e_2,e_4)$ contains four rays of total
degree $75+62+62+9=208$. Its subsheaf has slope $104>100$, proving
that $T_X$ is unstable.
\end{example}
